\documentclass[11pt]{article}
\usepackage[a4paper,margin=1in]{geometry}
\usepackage{amsmath,amsthm,amssymb}
\usepackage{mathtools}

\newtheorem{theorem}{Theorem}[section]
\newtheorem{proposition}[theorem]{Proposition}
\newtheorem{lemma}[theorem]{Lemma}
\newtheorem{corollary}[theorem]{Corollary}
\theoremstyle{definition}
\newtheorem{definition}[theorem]{Definition}
\newtheorem{remark}[theorem]{Remark}

\newcommand{\R}{\mathbb{R}}
\newcommand{\Fr}{\mathcal{A}}
\newcommand{\Om}{\Omega}
\newcommand{\eps}{\varepsilon}
\newcommand{\rstar}{\rho_{*}}
\newcommand{\rdel}{\rho_\delta}
\newcommand{\rhad}{\widehat\rho}
\newcommand{\X}{\mathcal X}

\title{Good-Endpoint Variant of the Weighted Metric Trace\\
       (Route~$\delta$, Plan~2, Building Block \textsc{GoodEndpointTrace})}
\author{}
\date{}

\begin{document}
\maketitle

\begin{abstract}
We strengthen the weighted metric trace of \cite{WeightedMetricTrace} by
producing a Borel level $\rhad\in G_I\cap G_\tau$ within
$O(\delta_T)$ of $\rdel=1-\kappa\sqrt{\delta_T}$ at which
$d_\X(F_{\rhad},B_1)^2\le C\delta_T$.  Here $G_I$ is the
Fusco--Julin good set $\{D_I\le\eta_I\}$ and
$G_\tau=\{w(\rho)\ge\tau_0\}$ is the profile-gap good set.  The point is
that the boundary-layer transfer of \cite{BoundaryLayerTransfer} only
requires the metric trace at \emph{some} level $\rhad$ with
$\rdel-\rhad=O(\delta_T)$: the collar volume
$|\Om\setminus E_{\rhad}|\le \omega_n n(\kappa\sqrt{\delta_T}+O(\delta_T))$
still has order $\sqrt{\delta_T}$, so squaring preserves the sharp
exponent.  This removes the need for $\rhad$ to coincide with the exact
endpoint $\rdel$.  Honest status of the bad-level kinetic hypotheses is
recorded in \S\ref{sec:bad-status}: both bad-$I$ action
\eqref{eq:bad-I-action} and bad-$\tau$ action
\eqref{eq:bad-tau-action} remain load-bearing for the underlying trace
theorem we invoke, and this variant does not by itself eliminate either.
\end{abstract}

\section{Notation and standing assumptions}\label{sec:notation}

We follow \cite{WeightedMetricTrace}.  Fix $n\ge 2$, $R\ge 2$,
$\rstar\in(0,1)$, $\kappa=\kappa(n,\rstar)>0$.  Let $\Om\subset B_R$ be
open with $|\Om|=\omega_n$, $u=u_\Om$ the variational torsion function,
$\delta:=\delta_T(\Om)$, $w(\rho)=-t'(\rho)$,
$d\mu(\rho)=w(\rho)d\rho$, $F_\rho:=\rho^{-1}E_\rho$.
The smallness threshold from
\cite[(eq:delta-smallness)]{WeightedMetricTrace} gives
\[
   \delta\le \delta_0:=\left(\frac{1-\rstar}{2\kappa}\right)^2,
   \qquad
   \rdel:=1-\kappa\sqrt\delta\ge\frac{1+\rstar}{2},
   \qquad
   L^*:=\frac{\rdel-\rstar}{4}\ge\frac{1-\rstar}{8}.
\]
Set $J^*:=[\rdel-L^*,\rdel]$, the order-one trace window.  Define
\begin{align*}
   G_I &:= \{\rho\in[\rstar,\rdel]:D_I(t(\rho))\le\eta_I\},
   & B_I &:= [\rstar,\rdel]\setminus G_I,\\
   G_\tau &:= \{\rho\in[\rstar,\rdel]:w(\rho)\ge\tau_0\},
   & B_\tau &:= [\rstar,\rdel]\setminus G_\tau,
   & \tau_0&:=\rstar/(2n),
\end{align*}
where $\eta_I=\eta_I(n,R,\rstar)>0$ is fixed in
\cite[Hyp.~3.5]{WeightedMetricTrace}.

\section{The good-endpoint trace}\label{sec:statement}

\begin{theorem}[Good-endpoint weighted metric trace]\label{thm:good-end}
There exist constants $C=C(n,R,\rstar,\kappa)>0$,
$C_1=C_1(n,R,\rstar,\kappa)>0$, and a smallness threshold
$\delta_0'(n,R,\rstar,\kappa)>0$ such that, under the conditional
hypotheses of \cite[Hyp.~3.5--3.7]{WeightedMetricTrace} (i.e.\ the
inputs that yield Theorem~3.3 of that block), the following holds.

For every open $\Om\subset B_R$ with $|\Om|=\omega_n$ and
$\delta=\delta_T(\Om)\le\delta_0'$, there exists a Borel level
\begin{equation}\label{eq:rhat-range}
   \rhad\;\in\;[\rdel-C_1\delta,\rdel]\cap G_I\cap G_\tau
\end{equation}
such that
\begin{equation}\label{eq:rhat-trace}
   d_\X(F_{\rhad},B_1)^{2}\;\le\;C(n,R,\rstar,\kappa)\,\delta.
\end{equation}
\end{theorem}

The proof has two steps: a measure-theoretic non-emptiness statement
for the candidate window, and a direct application of the existing
trace-uniform theorem of \cite{WeightedMetricTrace}.

\section{Bad-level measure bounds in Lebesgue form}\label{sec:bad-measure}

\begin{lemma}[Bad-set Lebesgue measures]\label{lem:bad-leb}
There is $C_{\rm bad}=C_{\rm bad}(n,R,\rstar)>0$ such that
\begin{align}
   \mathcal L^1(B_\tau)
      &\;\le\; C_{\rm bad}\,\delta,
      \label{eq:Btau-leb}\\
   \mathcal L^1(B_I\cap G_\tau)
      &\;\le\; C_{\rm bad}\,\delta.
      \label{eq:BI-Gtau-leb}
\end{align}
\end{lemma}

\begin{proof}
The first bound is exactly the profile-gap input
\cite[Hyp.~3.7, eq.\ (profile-gap)]{WeightedMetricTrace}.

For the second, the level-set deficit identity
\cite[(DH-DI-int)]{WeightedMetricTrace} combined with the definition of
$G_I$ gives
\[
   \mu(B_I)\;\le\;\eta_I^{-1}
      \int_{\rstar}^{\rdel}D_I(t(\rho))\,d\mu(\rho)
   \;\le\; \frac{C(n,\rstar)}{\eta_I}\,\delta.
\]
On $B_I\cap G_\tau$, $w(\rho)\ge\tau_0=\rstar/(2n)$, so
$\mu(B_I\cap G_\tau)\ge \tau_0\,\mathcal L^1(B_I\cap G_\tau)$.
Hence
\[
   \mathcal L^1(B_I\cap G_\tau)
   \;\le\;\frac{\mu(B_I)}{\tau_0}
   \;\le\;\frac{2nC(n,\rstar)}{\rstar\,\eta_I}\,\delta.
\]
Taking the maximum of the two constants gives $C_{\rm bad}$.
\end{proof}

\begin{lemma}[Non-empty good window near $\rdel$]\label{lem:window-nonempty}
Let $C_1:=3C_{\rm bad}$ with $C_{\rm bad}$ as in
Lemma~\ref{lem:bad-leb}.  Assume
$\delta\le\delta_0'(n,R,\rstar,\kappa):=\min\{\delta_0,L^*/(2C_1)\}$.
Then
\begin{equation}\label{eq:nonempty}
   \mathcal L^1\bigl([\rdel-C_1\delta,\rdel]\cap G_I\cap G_\tau\bigr)
   \;\ge\; C_{\rm bad}\,\delta\;>\;0.
\end{equation}
In particular the set in \eqref{eq:rhat-range} is non-empty and has
positive Lebesgue measure.
\end{lemma}

\begin{proof}
Write $I_\delta:=[\rdel-C_1\delta,\rdel]$.  The condition
$\delta\le L^*/(2C_1)$ ensures $C_1\delta\le L^*/2$, so
$I_\delta\subset J^*\subset[\rstar,\rdel]$, in particular
$\mathcal L^1(I_\delta)=C_1\delta$.
By Lemma~\ref{lem:bad-leb},
\[
   \mathcal L^1(I_\delta\cap B_\tau)\le C_{\rm bad}\delta,
   \qquad
   \mathcal L^1(I_\delta\cap B_I\cap G_\tau)
      \le \mathcal L^1(B_I\cap G_\tau)\le C_{\rm bad}\delta.
\]
The complement satisfies
\[
   I_\delta\cap G_I\cap G_\tau
   \;=\; I_\delta\setminus
      \bigl(B_\tau\cup (B_I\cap G_\tau)\bigr),
\]
and the two excluded sets are disjoint (the first sits in $B_\tau$, the
second in $G_\tau$).  Hence
\[
   \mathcal L^1(I_\delta\cap G_I\cap G_\tau)
   \;\ge\; C_1\delta - 2C_{\rm bad}\delta
   \;=\; (3C_{\rm bad}-2C_{\rm bad})\delta
   \;=\; C_{\rm bad}\delta\;>\;0.
\]
\end{proof}

\begin{remark}[Borel measurability]\label{rem:borel}
The set $G_I$ is Borel because $\rho\mapsto D_I(t(\rho))$ is Borel
(coarea on a Sobolev function); $G_\tau$ is Borel because $w=-t'$ is
Borel.  Hence the candidate set in \eqref{eq:rhat-range} is Borel, and
any specific Borel selector (e.g.\ the essential infimum of the set,
which is Borel-measurable because the set has positive measure)
provides a Borel-measurable choice of $\rhad$.
\end{remark}

\section{Proof of Theorem~\ref{thm:good-end}}\label{sec:proof}

\begin{proof}
Let $\delta\le\delta_0'$.  By Lemma~\ref{lem:window-nonempty}, the
candidate set in \eqref{eq:rhat-range} has positive Lebesgue measure;
choose any
$\rhad\in [\rdel-C_1\delta,\rdel]\cap G_I\cap G_\tau$ (Borel selection
exists by Remark~\ref{rem:borel}; the conclusion below does not depend
on the selection rule).
Since $C_1\delta\le L^*/2$, $\rhad\in J^*=[\rdel-L^*,\rdel]$.

The trace-uniform theorem
\cite[Thm.~3.3, eq.\ (trace-uniform)]{WeightedMetricTrace} states
\[
   \sup_{\rho\in J^*}d_\X(F_\rho,B_1)^2
   \;\le\; C_*(n,R,\rstar,\kappa)\,\delta.
\]
Specialising to $\rho=\rhad\in J^*$ gives
\eqref{eq:rhat-trace} with $C:=C_*$.
\end{proof}

\section{Verification: downstream boundary-layer transfer absorbs
$\rhad$}\label{sec:downstream}

The boundary-layer transfer of \cite{BoundaryLayerTransfer} is stated
at $\rdel$ but only uses the two ingredients
\eqref{eq:rhat-trace} and the collar volume bound.  We verify both
survive the replacement $\rdel\rightsquigarrow\rhad$.

\begin{lemma}[Collar volume at $\rhad$]\label{lem:collar-good}
For $\rhad$ as in Theorem~\ref{thm:good-end},
\begin{equation}\label{eq:collar-good}
   |\Om\setminus E_{\rhad}|\;\le\;\omega_n n\,(1-\rhad)
   \;\le\;\omega_n n\bigl(\kappa\sqrt\delta+C_1\delta\bigr)
   \;\le\;C_{\rm lay}'\sqrt\delta,
\end{equation}
where $C_{\rm lay}':=\omega_n n(\kappa+C_1\sqrt{\delta_0'})$.
\end{lemma}

\begin{proof}
Since $E_{\rhad}\subset\Om$ and $|E_{\rhad}|=\omega_n\rhad^{\,n}$,
$|\Om\setminus E_{\rhad}|=\omega_n(1-\rhad^{\,n})$.  The mean-value
theorem on $\rho\mapsto 1-\rho^n$ over $[\rhad,1]\subset[\rstar,1]$
gives $1-\rhad^{\,n}\le n(1-\rhad)$.  From \eqref{eq:rhat-range},
\[
   1-\rhad\;=\;(1-\rdel)+(\rdel-\rhad)\;\le\;\kappa\sqrt\delta+C_1\delta
   \;\le\;(\kappa+C_1\sqrt{\delta_0'})\sqrt\delta.
\]
\end{proof}

\begin{corollary}[Bounded sharp Saint--Venant via $\rhad$]\label{cor:SV-good}
Under the hypotheses of Theorem~\ref{thm:good-end}, there is
$C_{\rm SV}=C_{\rm SV}(n,R,\rstar,\kappa)>0$ such that
\begin{equation}\label{eq:SV-good}
   \Fr(\Om)^{2}\le C_{\rm SV}\,\delta.
\end{equation}
\end{corollary}

\begin{proof}
By the superlevel transfer lemma
\cite[Lemma~4.2]{BoundaryLayerTransfer} applied with $E_t=E_{\rhad}$
(valid because $E_{\rhad}\subset\Om$ and the lemma is purely
combinatorial in the volume defect),
\[
   \Fr(\Om)\;\le\;\Fr(E_{\rhad})+\frac{2}{\omega_n}|\Om\setminus E_{\rhad}|.
\]
By \eqref{eq:rhat-trace} and \cite[Rem.~2.4]{BoundaryLayerTransfer},
$\Fr(E_{\rhad})\le \omega_n^{-1}d_\X(F_{\rhad},B_1)$, so
$\Fr(E_{\rhad})^2\le \omega_n^{-2}C\,\delta$.  By
Lemma~\ref{lem:collar-good},
$|\Om\setminus E_{\rhad}|^2\le (C_{\rm lay}')^2\delta$.
Squaring with $(a+b)^2\le 2a^2+2b^2$,
\[
   \Fr(\Om)^2
   \;\le\; \frac{2C}{\omega_n^{2}}\delta
      + \frac{8(C_{\rm lay}')^2}{\omega_n^{2}}\,\delta
   \;=:\;C_{\rm SV}\,\delta.
\]
The exponent~$2$ is preserved exactly because the $O(\delta)$
contribution from $\rdel-\rhad$ is dominated by the
$O(\sqrt\delta)$ collar coming from $1-\rdel=\kappa\sqrt\delta$; the
cross term $O(\delta\cdot\sqrt\delta)=O(\delta^{3/2})$ is absorbed by
$O(\delta)$ for $\delta\le 1$.
\end{proof}

\section{Status of bad-level kinetic hypotheses}\label{sec:bad-status}

This section is the honest analysis the task brief asks for.  We
record exactly which of the two bad-level kinetic action bounds of
\cite[Hyp.~3.6]{WeightedMetricTrace}, namely
\begin{align}
   \int_{B_I}|\dot F_\rho|_\X^2\,d\mu(\rho)
      &\le C\delta,\tag{$K_I$}\label{eq:bad-I-action}\\
   \int_{B_\tau}|\dot F_\rho|_\X^2\,d\rho
      &\le C\delta,\tag{$K_\tau$}\label{eq:bad-tau-action}
\end{align}
are still load-bearing under the proof above.

\paragraph{What the proof above actually uses.}
The proof of Theorem~\ref{thm:good-end} invokes the trace-uniform
estimate \cite[Thm.~3.3]{WeightedMetricTrace} as a black box, hence it
inherits \emph{all} of that theorem's dependencies.  In the chain of
\cite{WeightedMetricTrace}:
\begin{itemize}
\item the integrated action bound (its Thm.~3.1) uses
\eqref{eq:bad-I-action} on $B_I$;
\item the Lebesgue kinetic bound (its Thm.~3.2) uses the integrated
action bound on $G_\tau$ plus \eqref{eq:bad-tau-action} on $B_\tau$;
\item the trace-uniform statement (its Thm.~3.3) uses both of the
above as (H1), (H4) in its abstract trace lemma.
\end{itemize}
Thus, taken at face value, the proof above does \emph{not} remove
either \eqref{eq:bad-I-action} or \eqref{eq:bad-tau-action}.

\paragraph{Can a refined proof remove \eqref{eq:bad-tau-action}?}
A potential refinement, sketched in the brief, is to repeat the
trace argument starting from a Markov point
$\rho_0\in J^*\cap G_\tau$ extracted from the integrated action
\eqref{eq:bad-I-action}-only bound, and then propagate to
$\rhad\in J^*\cap G_I\cap G_\tau$ by kinetic Cauchy--Schwarz on
$[\rhad,\rho_0]$.  The honest obstruction to dropping
\eqref{eq:bad-tau-action} is:
\begin{quote}
The interval $[\rhad,\rho_0]$ may intersect $B_\tau$, and the
Cauchy--Schwarz step needs
$\int_{[\rhad,\rho_0]}|\dot F_\rho|_\X^2\,d\rho<\infty$ with the
right $\delta$-scaling.  On $[\rhad,\rho_0]\cap G_\tau$ the bound
follows from $w\ge\tau_0$ and the integrated $d\mu$-action; on
$[\rhad,\rho_0]\cap B_\tau$ \emph{some} $\delta$-scaling input is
needed.
\end{quote}
The bound \(\mathcal L^1(B_\tau)\le C_{\rm bad}\delta\) controls only
the measure of $B_\tau$, not the $L^2(d\rho)$ integral of
$|\dot F_\rho|_\X$ over $B_\tau$.  Without \eqref{eq:bad-tau-action} or
a substitute (e.g.\ a pointwise Lipschitz bound on $\rho\mapsto F_\rho$
in $\X$, which is exactly the discredited A0), this Cauchy--Schwarz
step is not closed.  Hence the present good-endpoint variant does not
remove \eqref{eq:bad-tau-action} on its own.

\paragraph{Can a refined proof remove \eqref{eq:bad-I-action}?}
The bad-$I$ action input is used to bound the integrated action
$\int_{[\rstar,\rdel]}(d_\X^2+|\dot F|_\X^2)\,d\mu\le C\delta$ on $B_I$,
which then feeds Markov to find $\rho_0\in J^*\cap G_\tau$.  An
alternative route is to skip the integrated action entirely and
extract $\rho_0$ from the pointwise good-level estimate
$d_\X(F_\rho,B_1)^2\le C\,D_I(t(\rho))$ valid on $G_I$
(\cite[Hyp.~3.5]{WeightedMetricTrace}) combined with the integrated
budget $\int D_I\,d\mu\le C\delta$.  By Markov on
$d\mu|_{G_I\cap J^*}$ one obtains $\rho_0\in G_I\cap J^*$ with
$d_\X(F_{\rho_0},B_1)^2\le K_1\delta$ \emph{without using}
\eqref{eq:bad-I-action}.  The catch is that one then still has to
propagate this distance bound from $\rho_0$ to $\rhad$, which puts us
back at the same Cauchy--Schwarz obstruction on
$[\rhad,\rho_0]\cap B_\tau$.

In summary:
\begin{itemize}
\item Replacing the exact endpoint $\rdel$ by $\rhad\in G_I\cap G_\tau$
near $\rdel$ is a clean structural simplification of the
downstream interface (Corollary~\ref{cor:SV-good}).
\item It does \emph{not} by itself eliminate either bad-level kinetic
hypothesis.  Both \eqref{eq:bad-I-action} and
\eqref{eq:bad-tau-action} remain load-bearing for the trace-uniform
theorem that we invoke.
\item A genuine elimination of \eqref{eq:bad-tau-action} would require
a non-trivial substitute for the Cauchy--Schwarz transport across
$B_\tau$ (e.g., a quantitative AC bound for
$\rho\mapsto F_\rho$ at bad-density levels), which is exactly the
discredited A0; we do not provide one here.
\end{itemize}

\section{Conditionality summary}\label{sec:cond}

Theorem~\ref{thm:good-end} is conditional on the same external
ingredients as \cite[Thm.~3.3]{WeightedMetricTrace}:
\begin{itemize}
\item the level-set deficit identity \cite{LevelSetIdentity},
\item the conditional same-centre Fusco--Julin trace package
\cite[Hyp.~3.5]{WeightedMetricTrace},
\item the conditional good-level metric upper gradient
\cite[Hyp.~3.6]{WeightedMetricTrace},
\item the bad-level action inputs \eqref{eq:bad-I-action} and
\eqref{eq:bad-tau-action},
\item the profile-gap/Talenti density input
\cite[Hyp.~3.7]{WeightedMetricTrace}.
\end{itemize}
No new conditional input is introduced.  The variant produces a Borel
$\rhad$ in $G_I\cap G_\tau\cap[\rdel-C_1\delta,\rdel]$, which the
boundary-layer transfer of \cite{BoundaryLayerTransfer} absorbs without
loss of the sharp exponent~$2$, as shown in Corollary~\ref{cor:SV-good}.

\begin{thebibliography}{9}

\bibitem{WeightedMetricTrace}
Plan~2 building block \textsc{WeightedMetricTrace}.
Integrated action, kinetic bound, and weighted metric trace at
$\rdel$; see in particular Theorems~3.1--3.3, Hypotheses~3.5--3.7, and
the abstract trace Theorem~4.1.

\bibitem{BoundaryLayerTransfer}
Plan~2 building block \textsc{BoundaryLayerTransfer}.
Endpoint metric trace to bounded sharp Saint--Venant via superlevel
transfer; Theorem~3.1 (transfer at the level $\rhad$) and
Theorem~4.1 (small-$\delta_T$ Saint--Venant).

\bibitem{LevelSetIdentity}
Plan~2 building block \textsc{LevelSetIdentity}.
Exact level-set deficit identity
$\frac{2\delta_T}{|\Om|}=\frac{1}{|\Om|}\int(D_H+D_I)/(n^2\omega_n^{2/n}m^{1-2/n})\,dt$
and the profile-gap measure estimate
$\mathcal L^1(B_\tau)\le C\delta$.

\end{thebibliography}

\end{document}
